\scanpage{107}
\textcircled{2} \uline{Cas \deleted{\unsure{du}} lisse non quasi-projectif.}

Si $X$ est séparé, utilisant une compactification de Nagata, et résolution des singularités \inserted{comme dessus}, on a les \deleted{\unsure{ms}} résultats que ci-dessus, en notant que la résolution des singularités permet de déduire les \unsure{conj.|Conj.} de Weil pour des var. propres et lisses du cas projectif et lisse.

Si $X$ non séparé, $H^i_!(X)$ n'est plus défini. Par contre $H^i(X)$ est défini. On a, avec un recouvt affine
\[
 H^*(X)\Longleftarrow E^{pq}_2=H^p(\sigma\longleftarrow H^q(U_\sigma))
\]
Ici $E^{pq}_2$ est effectif de poids $\ge q$, mais \deleted{\unsure{on déduit peut-être ainsi|non dégénéré}} on peut en conclure que l'effectivité de $H^i(X)$, et le fait que ces poids sont compris entre $0$ et $2n$ si $i\ge n$, et entre $0$ et $2i$ si $i\le n$ ; en tous cas, ces poids sont $\le 2i$, \deleted{\unsure{enfin, le ni}}. Enfin, on trouve que le niveau est $\le i$ \deleted{\unsure{et sera donc resp. et $\le 2n-i$} (il suffit de regarder pour $i>n$)}

\margin{\uline{Question} le niveau est-il $\le 2n-i$ ? \deleted{\unsure{pour $X$ non séparé}}}
\margin{\unsure{ici|où|si} $n=\dim X$}

\textcircled{3} \uline{Cas non lisse} : $H^i_!(X)$, ($X$ séparé). Je dis qu'il est encore vrai qu'il est \uline{effectif} et de poids $\le i$, de niveau $\le\operatorname{Inf}(i,2n-i)$ ! Soit $U$ un ouvert lisse dense, $Y$ son complémentaire.

Donc
\[
 \to H^i_!(U,\unsure{\mathbb Q_\ell|{}})\to H^i_!(X)\to H^i_!(Y)\to
\]
et $H^i_!(U)$\deleted{\unsure{est}} est effectif de poids $\le i$ \inserted{de niveau $\le i$}, et $H^i_!(Y)$ aussi par hypothèse de récurrence sur la dimension, donc $H^i_!(X)$ est effectif de poids $\le i$, de niveau $\le i$. De \unsure{m} $H^i_!(U)$ est de niveau $\le 2n-i$, et $H^i_!(Y)$ de niveau $\le 2(n-1)-i\le 2n-i$, on gagne.

\textcircled{4} \uline{Cas non lisse} : $H^i(X)$ \deleted{\unsure{($X$ séparé)} \unsure{?}}. Je dis qu'il est effectif, à poids $\le\underline{\underline{2i}}$, à niveau \deleted{\unsure{$\le\operatorname{Inf}(i,n)$ où $n=\dim X$}} \inserted{$\le i$ et $\le n$}. \inserted{Par la suite spectrale de Leray, on est ramené au cas où $X$ affine.} Récurrence sur $i$, trivial si $i\le 0$. On utilise résolution $X'\to X$, et suite spectrale de descente
\[
 H^*(X)\Longleftarrow E_2^{pq}=H^p(s\longmapsto H^q((X'_{/X})^s))
\]
\margin{poids $\le 2\operatorname{Inf}(i,n)$}
\margin{niveau $\le\operatorname{Inf}(i,n)$}
\deleted{\unsure{si}} Par hyp. de récurrence, \deleted{\unsure{pour}} \inserted{si $q<i$}, $E_2^{pq}$ est effectif
% archive sheet number: 106
